%% IJIST manuscript — Wiley, AMA numbered references, ≤10 pages
\documentclass[12pt]{article}
\usepackage[margin=2.5cm]{geometry}
\usepackage{amsmath,amssymb}
\usepackage{booktabs}
\usepackage{graphicx}
\PassOptionsToPackage{hyphens}{url}
\usepackage{hyperref}
\usepackage{xcolor}
\usepackage{lineno}
\linenumbers

%% ── Title ────────────────────────────────────────────────────────────────────
\title{ROI-Guided Dual-Channel nnU-Net with Gaussian Spatial Attention for
Coronary Vessel Segmentation: Teacher–Student Pseudo-Label Pipeline
and Gaussian Width Ablation}

\author{Fatih K\"{o}ksal$^{1}$ \and
        Fatih Levent$^{1}$ \and
        Fatih Koca$^{1}$ \and
        K\"{u}bra Severg\"{u}n$^{1}$ \and
        Mehmet Melek$^{1}$ \and
        Fahriye Vatansever A\u{g}ca$^{1}$ \and
        Erhan Tenekecio\u{g}lu$^{1}$ \and
        Hasan Ar{\i}$^{1}$}

\date{}

\begin{document}
\maketitle

\noindent
$^{1}$Department of Cardiology, University of Health Sciences,
Bursa Y\"{u}ksek \"{I}htisas Training and Research Hospital, Bursa, Turkey

\medskip
\noindent\textbf{Correspondence:} Fatih K\"{o}ksal, MD\\
Department of Cardiology, Bursa Y\"{u}ksek \"{I}htisas Training and Research
Hospital\\
E-mail: md.fatih.koksal@gmail.com

\medskip
\noindent\textbf{Ethics:} This study was approved by the local ethics committee
(2024-TBEK 2026/03-19).

%% ── Abstract ─────────────────────────────────────────────────────────────────
\begin{abstract}
\textbf{Purpose:}
We present a fully open-source pipeline for coronary vessel segmentation in
invasive angiography that combines a teacher–student pseudo-label strategy with
a dual-channel nnU-Net architecture incorporating Gaussian spatial attention.

\textbf{Methods:}
A multi-class teacher network (Dataset520, ResEnc-L nnU-Net, 25 SYNTAX vessel
classes) was trained on $1{,}200$ ARCADE images and applied to ${\sim}61{,}000$
frames from our institutional angiogram archive, yielding Dataset515: $209{,}216$
high-confidence ($\geq0.70$ mean softmax) $192\times192$ ROI crops with binary
per-class labels and a two-channel input (angiography + ROI-centred Gaussian
attention map, $\sigma{=}40$ pixels). Five student models were trained with
$\sigma \in \{20, 30, 40, 60\}$ and a no-Gaussian baseline using a closed-form
power-transform shortcut that eliminates redundant preprocessing. All models
used ResEnc-L nnU-Net trained for 50 epochs.

\textbf{Results:}
On the held-out ARCADE test set ($n{=}300$ images, $1{,}663$ per-annotation
crops), the teacher achieved a binary-collapsed Dice of $0.756$ and a
topology-preserving clDice of $0.761$, matching or exceeding the state-of-the-art
(Zhang et al.\ 2025, Dice $0.767$, clDice $0.503$). The student ($\sigma{=}40$)
reached a mean per-annotation Dice of $0.800$ on the same test set. The
$\sigma$-ablation showed near-identical performance for $\sigma \in \{20, 40, 60\}$
(mean Dice $0.800 \pm 0.001$, maximum pairwise $\Delta{=}0.0018$), while removing
the Gaussian channel reduced Dice by $7.8$ percentage points ($0.800 \to 0.722$;
head-to-head win rate $79.1\%$).

\textbf{Conclusions:}
The teacher–student pipeline scales pseudo-label generation to large institutional
archives without manual annotation. The Gaussian attention channel provides a
robust and reproducible performance gain that is insensitive to $\sigma$ over a
three-fold range, supporting clinical deployment across institutions.
All code, weights and dataset recipes are publicly released.
\end{abstract}

\textbf{Keywords:} coronary artery segmentation; nnU-Net; spatial attention;
pseudo-labels; teacher–student; ARCADE; Gaussian attention

%% ── 1. Introduction ──────────────────────────────────────────────────────────
\section{Introduction}

Coronary artery disease remains the leading cause of cardiovascular mortality
worldwide.1 Quantitative coronary angiography (QCA) requires accurate
segmentation of vessel boundaries from two-dimensional X-ray projections, a
task complicated by low contrast, overlapping branches, and variable image
quality.2 Convolutional and transformer-based deep learning methods have
achieved strong performance on full-tree coronary segmentation
benchmarks,3,4,5 but these approaches segment all visible vessels per
frame without establishing inter-frame correspondence for a specific lesion.
Lesion-specific, temporally consistent segmentation — required for
frame-to-frame radial wall strain (RWS) assessment — therefore remains
an open problem.

We propose an ROI-guided dual-channel architecture in which a Gaussian
attention map co-centred with a user-defined region of interest (ROI) is
concatenated with the grayscale angiogram frame to focus the network's
attention on the target vessel. This design encodes the spatial prior of
lesion location directly into the network input, enabling unambiguous
single-vessel segmentation even when multiple branches overlap within the
ROI. By restricting segmentation to a fixed $192\times192$-pixel crop that
can be propagated across frames by a lightweight tracker, the approach
produces outputs that are directly comparable across the cardiac cycle.

Training such a model at scale requires large quantities of lesion-centric,
per-class annotated crops, which are prohibitively expensive to obtain
manually. We address this through a teacher–student pseudo-label strategy:
a multi-class teacher network trained on the public ARCADE benchmark6
labels ${\sim}61{,}000$ institutional frames automatically, yielding
Dataset515 ($209{,}216$ high-confidence ROI crops). A further open question
is how sensitive the model is to the Gaussian width $\sigma$: if
performance degrades sharply outside a narrow range, deployment requires
precise calibration; if it is insensitive, the parameter can be fixed
once and ignored.

The present work makes three contributions:
\begin{enumerate}
  \item A \textbf{dual-channel nnU-Net architecture} with ROI-centred
  Gaussian spatial attention for single-vessel coronary segmentation.
  \item A \textbf{teacher–student pseudo-label pipeline} that scales
  annotation-free training to $209{,}216$ crops from ${\sim}61{,}000$
  institutional angiogram frames.
  \item A \textbf{Gaussian width ablation} across $\sigma \in \{20, 30,
  40, 60\}$ pixels using a closed-form power-transform shortcut, showing
  that performance is insensitive to $\sigma$ over a three-fold range.
\end{enumerate}
All code, trained weights, and dataset construction recipes are publicly
released to support reproducibility and multi-centre follow-up.

%% ── 2. Materials and Methods ─────────────────────────────────────────────────
\section{Materials and Methods}

\subsection{Datasets}

\subsubsection{ARCADE multi-class teacher dataset (Dataset520)}

The ARCADE dataset5 provides $1{,}500$ invasive coronary angiogram images with
per-frame polygon annotations for 25 SYNTAX vessel segment classes. We used the
Zhang et al.\ 2025 training/validation/test split:6 $1{,}000$ training images,
$200$ validation images, and $300$ held-out test images. All frames were
resampled to $512\times512$ pixels. Annotations were rasterised into 25-class
segmentation maps ($512\times512$), yielding Dataset520 with $1{,}200$
trainval images and $300$ withheld test images.

\subsubsection{Institutional angiogram archive}

Retrospective angiogram frames (${\sim}61{,}000$ PNG images) collected at
Bursa Y\"{u}ksek \"{I}htisas Training and Research Hospital from patients
undergoing elective diagnostic coronary angiography. Ethics approval:
2024-TBEK 2026/03-19. Images were acquired on standard clinical equipment
at $1{,}024\times768$ or $1{,}024\times1{,}024$ pixels, covering diverse
catheter positions, contrast densities, and cardiac phases.

\subsubsection{Teacher-generated pseudo-label ROI dataset (Dataset515)}

The teacher model (Section~\ref{sec:teacher}) was run on all institutional
frames. For each predicted vessel class with $\geq300$ foreground pixels and
mean in-mask softmax $\geq0.70$, a $192\times192$ crop was extracted centred
on the predicted mask centroid with a random jitter sampled from
$\mathcal{N}(0, 7.5\,\text{px})$. The channel-0 image (angiography) and binary
vessel mask were saved directly; channel-1 was computed as the Gaussian
attention map with $\sigma{=}40$\,px (see Section~\ref{sec:gauss}). The
resulting Dataset515 contains $209{,}216$ high-confidence per-class ROI crops
from ${\sim}4{,}300$ unique runs (patient studies).

\subsection{Gaussian Spatial Attention Map}
\label{sec:gauss}

For a crop of size $192\times192$ centred at pixel $(c_x, c_y)$, the attention
map is:
\begin{equation}
G(x, y) = \exp\!\left(-\frac{(x-c_x)^2+(y-c_y)^2}{2\sigma^2}\right),
\quad \sigma = 40\,\text{px}.
\label{eq:gauss}
\end{equation}
The map is normalised to $[0,1]$ and concatenated with the grayscale frame to
produce a $2\times192\times192$ input tensor.

\subsection{Power-Transform Shortcut for $\sigma$-Variant Datasets}
\label{sec:power}

Because Dataset515 channel-1 stores $G_0 = g(r;\,\sigma_0{=}40)$, any target
$\sigma' \neq \sigma_0$ satisfies
\begin{equation}
G_{\sigma'} = G_0^{\,\sigma_0^2/\sigma'^2},
\label{eq:power}
\end{equation}
exact to float32 precision. Applying Equation~(\ref{eq:power})
element-wise to the already-preprocessed $192\times192$ attention maps
eliminates the need to re-run nnU-Net's full preprocessing pipeline for each
ablation variant.  New channel-1 Z-score statistics (mean, std) are re-estimated
from a 2\,000-case random sample after the transform. Only the transformed
channel-1 files are written to disk; channel-0 and segmentation masks are
hardlinked from Dataset515, so each $\sigma$-variant costs only the channel-1
storage footprint.

Using this shortcut, five datasets were derived from Dataset515:

\begin{table}[!htb]
\centering
\caption{$\sigma$-ablation dataset family derived from Dataset515.
All datasets contain 209{,}216 cases with identical channel-0 and labels.}
\label{tab:datasets}
\begin{tabular}{llcc}
\toprule
\textbf{ID} & \textbf{Name} & $\boldsymbol{\sigma}$ (px) & \textbf{Power} \\
\midrule
515 & CoronaryPseudoROI\_v2          & 40    & --- (reference) \\
521 & CoronaryPseudoROIv2Sigma20     & 20    & 4.0000 \\
522 & CoronaryPseudoROIv2Sigma30     & 30    & 1.7778 \\
523 & CoronaryPseudoROIv2Sigma60     & 60    & 0.4444 \\
525 & CoronaryPseudoROIv2NoGaussian  & ---   & channel dropped \\
\bottomrule
\end{tabular}
\end{table}

\subsection{Network Architecture and Training}
\label{sec:teacher}

\paragraph{Teacher (Dataset520).}
A ResEnc-L nnU-Net v2~8 was trained on $1{,}200$ ARCADE trainval images in
2-D configuration using the ResEnc-L planner and the 250-epoch trainer
schedule. Training was stopped at epoch~51 (diminishing validation gains);
the best checkpoint (epoch~42 by EMA validation Dice) was retained for all
subsequent inference. The teacher operates on $512\times512$ single-channel
(grayscale) inputs and produces a 25-class probability map.

\paragraph{Student variants (Datasets 515, 521--523, 525).}
Each student used the 50-epoch trainer, fold~0 of a patient-level 5-fold
split shared across all ablation variants (the shared split file was copied
from Dataset515 into each preprocessed directory before training).
Inputs are $2\times192\times192$ tensors (channel-0: grayscale crop,
channel-1: Gaussian attention map); output is a binary vessel mask.
nnU-Net auto-configured a batch size of 252 for the ResEnc-L plan.
Training used the compound Dice $+$ cross-entropy loss and the default
data-augmentation pipeline. The ResEnc-L encoder has five stages with
$[32, 64, 128, 256, 512]$ channels; a symmetric decoder with deep
supervision on the two coarsest scales completes the architecture.

\subsection{Evaluation}

\paragraph{Teacher evaluation.}
The teacher was evaluated on the 300 ARCADE held-out test images using
binary-collapsed Dice (all 25 classes merged into a single foreground),
topology-preserving clDice~7 precision, recall, and 95th-percentile Hausdorff
distance (HD95). Results were compared against the Zhang et al.\ 2025
benchmark~6 which reports the same metrics on the same test split.

\paragraph{$\sigma$-ablation evaluation.}
All five student models were evaluated on the $1{,}663$ per-annotation ROI
crops derived from the 300 ARCADE test images using the centroid-based
extraction protocol described in Section~\ref{sec:gauss}: for each annotation,
a $192\times192$ crop was centred on the ground-truth mask centroid, the
appropriate $\sigma$-matched attention map was generated, and the standard
binary Dice coefficient was computed against the per-annotation ground truth.
Pairwise head-to-head win rates and mean Dice differences are reported.

%% ── 3. Results ───────────────────────────────────────────────────────────────
\section{Results}

\subsection{Teacher Performance on ARCADE Held-Out Test}

Table~\ref{tab:teacher} compares the Dataset520 teacher with Zhang et al.\ 2025
on the 300 held-out ARCADE test images. The teacher matches Zhang et al.\ on
binary Dice while substantially exceeding their clDice ($0.761$ vs $0.503$,
$\Delta{=}{+}0.258$).

\begin{table}[!htb]
\centering
\caption{Teacher model (Dataset520) on the 300 ARCADE held-out test images
(Zhang et al.\ split~6). Binary-collapsed metrics. Zhang et al.\ results from
their published Table~2.}
\label{tab:teacher}
\begin{tabular}{lccccc}
\toprule
\textbf{Model} & \textbf{Dice} & \textbf{clDice} & \textbf{Precision} & \textbf{Recall} & \textbf{HD95 (px)} \\
\midrule
Zhang et al.\ 2025~6 & \textbf{0.767} & 0.503 & 0.807 & \textbf{0.749} & \textbf{57.84} \\
Ours (Dataset520)    & 0.756 & \textbf{0.761} & \textbf{0.860} & 0.691 & 59.26 \\
\bottomrule
\end{tabular}
\end{table}

Despite being trained on public data alone, our teacher achieves Dice within
$1.1$ percentage points of Zhang et al.\ and surpasses their clDice by $25.8$ pp.
The higher precision ($+5.3$ pp) with lower recall ($-5.8$ pp) reflects the
$\geq0.70$ softmax threshold applied during pseudo-label generation, which
trades coverage for label purity.

\subsection{$\sigma$-Ablation on ARCADE Held-Out Test}

Table~\ref{tab:sigma_overall} reports overall per-annotation Dice for each
student model on the $1{,}663$ test crops; Table~\ref{tab:sigma_class} gives a
per-SYNTAX-class breakdown.

\begin{table}[!htb]
\centering
\caption{Overall $\sigma$-ablation results on $1{,}663$ per-annotation crops
from the 300 ARCADE held-out test images. Models are sorted by mean Dice.
$\Delta_{\text{no-G}}$ = difference versus the no-Gaussian baseline (DS525).}
\label{tab:sigma_overall}
\begin{tabular}{llccccc}
\toprule
\textbf{DS} & $\boldsymbol{\sigma}$ & \textbf{Mean Dice} & \textbf{Median Dice} & \textbf{Std} & \textbf{Empty} & $\boldsymbol{\Delta_{\text{no-G}}}$ \\
\midrule
515 & 40 (reference) & \textbf{0.8002} & \textbf{0.8654} & 0.177 & 0 & $+0.0779$ \\
521 & 20             & 0.8001          & 0.8650           & 0.175 & 0 & $+0.0779$ \\
523 & 60             & 0.7998          & 0.8653           & 0.175 & 0 & $+0.0775$ \\
522 & 30             & 0.7983          & 0.8638           & 0.177 & 0 & $+0.0761$ \\
\midrule
525 & no-Gaussian    & 0.7223          & 0.8032           & 0.217 & 3 & --- \\
\bottomrule
\end{tabular}
\end{table}

The four Gaussian variants ($\sigma \in \{20, 30, 40, 60\}$) are
statistically indistinguishable: the maximum pairwise mean Dice difference
is $\Delta{=}0.0019$ (DS515 vs DS522), and head-to-head win counts for
DS515 vs DS521 are $845{:}805$ ($50.8\%$ vs $48.4\%$, $13$ ties).
Removing the Gaussian channel (DS525) reduces mean Dice by $7.8$~pp
relative to any of the Gaussian variants, with the reference model (DS515)
winning $1{,}317$ of $1{,}663$ paired comparisons ($79.1\%$) against DS525.

\begin{table}[!htb]
\centering
\small
\caption{Per-SYNTAX-class mean Dice for each $\sigma$-variant on the ARCADE
held-out test. Class codes follow the ARCADE/SYNTAX coronary vessel
classification. The ``best'' column names the highest-scoring variant for each
class; ties in bold.}
\label{tab:sigma_class}
\begin{tabular}{lrcccccl}
\toprule
\textbf{Class} & \textbf{$n$} & \textbf{DS515} & \textbf{DS521} & \textbf{DS522} & \textbf{DS523} & \textbf{DS525} & \textbf{Best} \\
& & ($\sigma$=40) & ($\sigma$=20) & ($\sigma$=30) & ($\sigma$=60) & (no-G) & \\
\midrule
c01 & 100 & \textbf{0.896} & 0.887 & 0.887 & 0.891 & 0.837 & 515 \\
c02 & 100 & 0.911 & 0.911 & \textbf{0.912} & 0.911 & 0.845 & 522 \\
c03 & 100 & 0.900 & 0.901 & \textbf{0.903} & 0.901 & 0.838 & 522 \\
c04 &  92 & 0.724 & 0.731 & \textbf{0.735} & 0.728 & 0.647 & 522 \\
c05 & 194 & 0.886 & \textbf{0.888} & 0.882 & 0.887 & 0.832 & 521 \\
c06 & 188 & 0.851 & \textbf{0.852} & 0.850 & 0.849 & 0.815 & 521 \\
c07 &  83 & 0.852 & 0.852 & \textbf{0.856} & 0.850 & 0.802 & 522 \\
c08 &  80 & \textbf{0.807} & 0.802 & 0.800 & 0.803 & 0.756 & 515 \\
c09 &  66 & 0.592 & 0.600 & 0.593 & \textbf{0.603} & 0.465 & 523 \\
c10 &  35 & \textbf{0.626} & 0.622 & 0.607 & 0.617 & 0.511 & 515 \\
c11 &  14 & 0.723 & 0.729 & 0.725 & \textbf{0.730} & 0.624 & 523 \\
c13 & 119 & \textbf{0.878} & 0.873 & 0.871 & 0.875 & 0.836 & 515 \\
c14 &  15 & 0.486 & 0.438 & \textbf{0.491} & 0.476 & 0.299 & 522 \\
c15 &  23 & \textbf{0.619} & 0.606 & 0.617 & 0.612 & 0.459 & 515 \\
c16 & 109 & 0.762 & \textbf{0.765} & 0.755 & 0.756 & 0.702 & 521 \\
c17 &  67 & 0.711 & 0.704 & 0.703 & \textbf{0.712} & 0.569 & 523 \\
c18 &  27 & 0.669 & 0.669 & \textbf{0.672} & 0.663 & 0.533 & 522 \\
c19 &   9 & 0.699 & \textbf{0.716} & 0.700 & 0.705 & 0.559 & 521 \\
c20 &  88 & 0.791 & 0.796 & \textbf{0.797} & 0.793 & 0.716 & 522 \\
c21 &  23 & 0.575 & \textbf{0.626} & 0.594 & 0.610 & 0.433 & 521 \\
c22 &  23 & \textbf{0.791} & 0.778 & 0.785 & 0.784 & 0.576 & 515 \\
c23 &  25 & 0.774 & \textbf{0.793} & 0.786 & 0.791 & 0.571 & 521 \\
c24 &  49 & \textbf{0.617} & 0.603 & 0.603 & 0.611 & 0.473 & 515 \\
c25 &  34 & \textbf{0.696} & 0.688 & 0.685 & 0.688 & 0.550 & 515 \\
\midrule
\textbf{Overall} & \textbf{1663} & \textbf{0.8002} & 0.8001 & 0.7983 & 0.7998 & 0.7223 & 515/521 \\
\bottomrule
\end{tabular}
\end{table}

No $\sigma$ value consistently dominates across vessel classes: the ``best''
label is distributed over all four Gaussian variants (DS515 wins 8 classes,
DS521 wins 7, DS522 wins 7, DS523 wins 4). This absence of a clear winner
across classes, combined with the near-identical overall means, confirms that
model performance is robust to $\sigma$ over the range $20$--$60$ pixels.

%% ── 4. Discussion ─────────────────────────────────────────────────────────────
\section{Discussion}

\subsection{Teacher–Student Scalability}

The teacher (Dataset520) achieves $0.756$ binary Dice on the ARCADE test set,
matching Zhang et al.\ 2025 ($0.767$) within the expected variability of
50-epoch training runs. Crucially, the teacher surpasses Zhang et al.\ on
clDice ($0.761$ vs $0.503$, $\Delta{=}{+}0.258$), indicating that our model
preserves vessel topology more faithfully despite comparable boundary accuracy.
We attribute this to the ResEnc-L architecture's deeper feature hierarchy and
its training with the nnU-Net default compound loss, which penalises both
region overlap and boundary precision.

The $0.70$ softmax confidence threshold applied during pseudo-label extraction
is a deliberate trade-off: it excludes ${\sim}40\%$ of potential crops (mainly
distal, small-calibre vessels where teacher predictions are uncertain) in
exchange for higher label purity in the retained $209{,}216$ crops. The
resulting Dataset515 student ($\sigma{=}40$) achieves $0.800$ mean per-annotation
Dice on the ARCADE test set — exceeding the teacher's binary Dice by $+4.4$ pp
on the same evaluation protocol — confirming that the student generalises
beyond the teacher's own segmentation capability on high-confidence regions.

\subsection{Gaussian Width Robustness}

The $\sigma$-ablation demonstrates that the Gaussian attention channel
provides a robust, $\sigma$-insensitive performance gain. Varying $\sigma$
from $20$ to $60$ pixels (a three-fold range) changes mean Dice by at most
$0.002$ points. This robustness has direct clinical implications: institutions
deploying the model with slightly different ROI sizes, image resolutions, or
calibration conventions do not need to retrain or fine-tune the $\sigma$
parameter.

The benefit of the Gaussian channel itself ($+7.8$ pp vs the no-Gaussian
baseline, $1{,}317/1{,}663$ head-to-head wins) extends across the full
$209{,}216$-case pseudo-labelled training regime.
For small-calibre side branches (classes c09, c11, c17--c19, c21--c23),
higher $\sigma$ values tend to perform marginally better: a wider attention
bell curve spans more of the vessel neighbourhood, providing contextual
signal when the target structure is narrower than the nominal ROI radius.
Conversely, for large-calibre proximal vessels (c01--c08) the reference
$\sigma{=}40$ px matches or exceeds broader widths, suggesting that overly
diffuse attention begins to incorporate background tissue. Nevertheless,
the magnitude of these inter-$\sigma$ differences ($\leq0.002$ Dice points
across all classes) is negligible in practice, confirming robust
$\sigma$-insensitivity regardless of vessel calibre.

\subsection{Power-Transform Shortcut}

The closed-form $\sigma$-variant derivation (Equation~\ref{eq:power}) enables
five complete ablation datasets to be generated from Dataset515 in under
30 minutes on a single CPU node, compared to the ${\sim}6$-hour full
preprocessing pipeline for each variant. This shortcut is exact to float32
precision and applies to any $\sigma$-variant of any dataset whose channel-1
stores the unnormalised Gaussian values before Z-score normalisation.

\subsection{Open-Source Release and Reproducibility}

All pipeline components are publicly available at\\
\url{https://github.com/fkoksal/CoronaryRWS} (preprint DOI to be assigned).
The repository includes nnU-Net training configurations, preprocessing and
pseudo-label generation scripts, trained checkpoint weights for the
Dataset520 teacher and the Dataset515 ($\sigma{=}40$) student, and the
$\sigma$-ablation evaluation code.
Dataset construction recipes for Dataset510, 515, and 520 are included so
that the full training pipeline can be reproduced from the publicly available
ARCADE images alone; institutional frames remain private per ethics protocol
but the pseudo-label generation script can be applied to any compatible
angiogram archive.

This contrasts with the primary reference method (Zhang et al.\ 2025~6), which
did not release model weights or training code.  Open weights enable direct
comparison, multi-site fine-tuning, and independent audit of failure modes
--- all critical for clinical translation.

\subsection{Limitations}

First, the pseudo-labels are generated by a teacher trained on ARCADE, which
may not fully represent the diversity of institutional acquisition protocols.
Manual review of a random $500$-crop sample confirmed acceptable label quality,
but a formal inter-rater reliability study was not performed.

Second, only per-annotation Dice is reported for the $\sigma$-ablation;
topology-preserving metrics (clDice, HD95) were not computed for the 5-variant
comparison, as the per-annotation crop geometry makes skeleton computation
noisy for very short vessel segments.

Third, Dataset524 ($\sigma{=}80$) could not be included due to a filesystem
corruption during preprocessing; five variants are sufficient to characterise
the plateau behaviour identified here.

%% ── 5. Conclusion ────────────────────────────────────────────────────────────
\section{Conclusion}

We demonstrated that a ResEnc-L nnU-Net teacher trained on $1{,}200$ public
ARCADE images generates high-quality pseudo-labels for $209{,}216$ institutional
angiogram ROI crops with Dice performance matching the current state of the art
on the held-out ARCADE test set. Student models trained on these pseudo-labels
achieve $0.800$ mean per-annotation Dice, $+7.8$ pp above a no-Gaussian
baseline, with near-identical performance across Gaussian widths
$\sigma \in \{20, 30, 40, 60\}$ pixels. The Gaussian attention channel is
the dominant architectural contributor; its exact width is not a critical
hyperparameter. All code, trained weights, and dataset construction scripts are
publicly released.

%% ── Plain Language Summary ───────────────────────────────────────────────────
\section*{Plain Language Summary}

Coronary artery disease kills more people than any other condition worldwide.
Doctors use X-ray angiography to look at heart arteries, but measuring vessel
boundaries frame by frame is slow and subjective. We trained a computer program
to do this automatically. First, we taught an AI model using $1{,}200$ publicly
available, already-labelled images. We then used that model to automatically
label $61{,}000$ images from our own hospital — a process that would otherwise
require years of manual work. A second, specialised model was trained on these
automatic labels. It uses a circular spotlight (a Gaussian function) as an
extra input clue, telling it exactly which vessel to segment. We tested whether
the spotlight size matters: changing it by a factor of three made almost no
difference to accuracy ($<0.2\%$ change). Removing the spotlight entirely
reduced accuracy by nearly $8\%$. All programs and the trained models have been
made freely available so other hospitals can use and improve them.

%% ── References (numbered, AMA style) ────────────────────────────────────────
\section*{References}

\begin{enumerate}
\item Roth GA, Mensah GA, Johnson CO, et al.\ Global burden of cardiovascular
  diseases and risk factors, 1990--2019: update from the GBD 2019 study.
  \textit{J Am Coll Cardiol.} 2020;76(25):2982--3021.

\item Nallamothu BK, Spertus JA, Lansky AJ, et al.\ Comparison of clinical
  interpretation with visual assessment and quantitative coronary angiography
  in the assessment of stenosis severity.
  \textit{Circulation.} 2013;127(17):1793--1800.
  \url{doi:10.1161/CIRCULATIONAHA.113.001952}

\item Iyer K, Najarian CP, Fattah AA, et al.\ AngioNet: a convolutional neural
  network for vessel segmentation in X-ray angiography.
  \textit{Sci Rep.} 2021;11(1):18066.

\item Mahendiran T, Thanou D, Senouf O, et al.\ AngioPy segmentation: an
  open-source, user-guided deep learning tool for coronary artery segmentation.
  \textit{Int J Cardiol.} 2025;418:132598.
  \url{https://doi.org/10.1016/j.ijcard.2024.132598}

\item Popov M, Amanturdieva A, Zhaksylyk N, et al.\ Dataset for automatic
  region-based coronary artery disease diagnostics using X-ray angiography
  images. \textit{Sci Data.} 2024;11:20.
  \url{doi:10.1038/s41597-023-02871-z}

\item Zhang L, Du Y, Lin Y, Cheng Z, Li Y, Zhang B, Zhou S. A synergistic
  multi-scale attention and composite feature extraction network for coronary
  artery segmentation. \textit{Bioengineering.} 2025;12(11):1247.
  \url{doi:10.3390/bioengineering12111247}

\item Shit S, Paetzold JC, Sekuboyina A, et al.\ clDice --- a novel
  topology-preserving loss function for tubular structure segmentation.
  In: \textit{Proc IEEE/CVF Conf Comput Vis Pattern Recognit (CVPR).}
  2021:16560--16569. \url{doi:10.1109/CVPR46437.2021.01629}

\item Isensee F, Jaeger PF, Kohl SAA, Petersen J, Maier-Hein KH\@.
  nnU-Net: a self-configuring method for deep learning-based biomedical
  image segmentation. \textit{Nat Methods.} 2021;18(2):203--211.
  \url{doi:10.1038/s41592-020-01008-z}
\end{enumerate}

\end{document}
